% ============================================================
%  sections/chapter_04.tex
%  Section 4: Norm Analysis
% ============================================================

\section{Norm Analysis}
\label{sec:norms}
The choice of reconstruction loss $\mathcal{L}$ is not merely a
computational convenience: it encodes an explicit assumption about the
noise distribution and determines which errors the approximation
prioritises minimising.
This section analyses the three losses introduced in
Section~\ref{subsec:losses} from statistical, geometric, and numerical
perspectives, and derives the implications for Tucker reconstruction
under each norm.

% ------------------------------------------------------------
\subsection{Frobenius Norm: Statistical and Geometric Properties}
\label{subsec:frob_analysis}
% ------------------------------------------------------------

\paragraph{Maximum likelihood interpretation.}
Suppose the observed tensor is
$\tensor{X} = \hat{\tensor{X}} + \tensor{E}$,
where $e_{i_1\cdots i_N} \overset{\text{i.i.d.}}{\sim} \mathcal{N}(0, \sigma^2)$.
The log-likelihood is
\begin{equation}
  \log p(\tensor{X} \mid \hat{\tensor{X}})
  \;=\; -\frac{1}{2\sigma^2} \frobn{\tensor{X} - \hat{\tensor{X}}}^2
        + \text{const},
  \label{eq:gauss_ll}
\end{equation}
so the maximum likelihood estimator of $\hat{\tensor{X}}$ coincides
exactly with the Frobenius minimiser.
Under this noise model, HOOI is optimal.

\paragraph{Unitary invariance and the ALS simplification.}
Unitary invariance~\eqref{eq:unitary_inv} implies that when the factor
matrices are orthonormal, the objective decomposes as
\begin{equation}
  \frobn{\tensor{X} - \hat{\tensor{X}}}^2
  \;=\; \frobn{\tensor{X}}^2 - \frobn{\tensor{G}}^2,
  \label{eq:frob_decomp}
\end{equation}
so \emph{minimising the reconstruction error is equivalent to maximising
the energy captured by the core tensor}.
This is the variational characterisation exploited by HOSVD and HOOI.

\paragraph{Sensitivity to outliers.}
The squared loss assigns weight proportional to $|r|$ to each residual
$r = x - \hat{x}$, meaning that a single large outlier can dominate the
objective and distort all factor matrices.
Formally, the \emph{influence function} of the Frobenius estimator is
unbounded~\cite{hampel1986robust}:
\begin{equation}
  \mathrm{IF}(r;\, \mathcal{L}_F) \;=\; r,
  \label{eq:if_frob}
\end{equation}
confirming zero robustness to gross errors.

% ------------------------------------------------------------
\subsection{\texorpdfstring{$\ell_1$}{L1} Norm: Robustness and
            Non-Smoothness}
\label{subsec:l1_analysis}
% ------------------------------------------------------------

\paragraph{Maximum likelihood interpretation.}
If $e_{i_1\cdots i_N} \overset{\text{i.i.d.}}{\sim}
\mathrm{Laplace}(0, b)$, the log-likelihood is
\begin{equation}
  \log p(\tensor{X} \mid \hat{\tensor{X}})
  \;=\; -\frac{1}{b}\onenorm{\tensor{X} - \hat{\tensor{X}}}
        + \text{const},
  \label{eq:laplace_ll}
\end{equation}
making the $\ell_1$ minimiser the maximum likelihood estimator under
Laplace noise.
The heavier tails of the Laplace distribution relative to the Gaussian
reflect the outlier tolerance built into $\mathcal{L}_1$.

\paragraph{Breakdown point.}
The \emph{breakdown point} $\varepsilon^*$ of an estimator is the fraction
of arbitrarily corrupted observations it can tolerate before becoming
arbitrarily biased.
The Frobenius estimator has $\varepsilon^* = 0$ (a single outlier suffices
to push the estimate arbitrarily far), whereas the $\ell_1$ estimator
achieves $\varepsilon^* = 0.5$ for location problems~\cite{maronna2019robust}.
In the Tucker setting, this translates to robustness against sparse
corruption of up to half the tensor entries.

\paragraph{Subdifferential and optimality condition.}
Because $\mathcal{L}_1$ is non-smooth, classical gradient conditions do
not apply.
The subdifferential of $|r|$ at $r$ is
\begin{equation}
  \partial|r|
  \;=\; \begin{cases}
          \{\sign(r)\} & r \neq 0, \\
          [-1,\, 1]    & r = 0,
        \end{cases}
  \label{eq:subgrad_l1}
\end{equation}
and the optimality condition for~\eqref{eq:general_opt} with
$\mathcal{L} = \mathcal{L}_1$ requires that zero lie in the
subdifferential of the objective with respect to each factor.
This is the condition that IRLS (Section~\ref{subsec:irls}) seeks to
satisfy at convergence.

\paragraph{Influence function.}
The influence function of the $\ell_1$ estimator is bounded:
\begin{equation}
  \mathrm{IF}(r;\, \mathcal{L}_1) \;=\; \sign(r),
  \label{eq:if_l1}
\end{equation}
confirming that a single outlier cannot push the estimate beyond a finite
neighbourhood of the truth.

% ------------------------------------------------------------
\subsection{Huber Loss: Smooth Interpolation}
\label{subsec:huber_analysis}
% ------------------------------------------------------------

\paragraph{Definition and properties.}
The Huber loss $\rho_\delta$~\eqref{eq:huber} satisfies:
\begin{itemize}
  \item $\rho_\delta \in C^1(\R)$: continuously differentiable everywhere,
        enabling gradient descent without subgradient complications.
  \item $\rho_\delta(r) = \tfrac{1}{2}r^2$ for $|r| \leq \delta$:
        quadratic near zero, consistent with Gaussian small errors.
  \item $\rho_\delta(r) \sim \delta|r|$ for $|r| \gg \delta$:
        linear in the tails, inheriting the robustness of $\ell_1$.
  \item $\rho_\delta \to \tfrac{1}{2}|\cdot|^2$ as $\delta \to \infty$
        and $\rho_\delta \to \delta|\cdot|$ as $\delta \to 0$:
        Huber interpolates between pure Frobenius and pure $\ell_1$.
\end{itemize}

\paragraph{Influence function.}
\begin{equation}
  \mathrm{IF}(r;\, \mathcal{L}_H)
  \;=\; \psi_\delta(r)
  \;=\; \begin{cases}
          r                    & |r| \leq \delta, \\
          \delta\,\sign(r) & |r| > \delta,
        \end{cases}
  \label{eq:if_huber}
\end{equation}
which is bounded (unlike Frobenius) yet smooth (unlike $\ell_1$).

\paragraph{Threshold selection.}
The parameter $\delta$ controls the boundary between the quadratic and
linear regimes.
A data-adaptive choice is
\begin{equation}
  \delta^*
  \;=\; c \cdot \widehat{\mathrm{MAD}}(\tensor{R}),
  \qquad
  \widehat{\mathrm{MAD}} = \mathrm{median}\!\left(|r_{i_1\cdots i_N}|
  \right),
  \label{eq:delta_mad}
\end{equation}
where $c \approx 1.345$ is the standard efficiency constant for
Gaussian noise~\cite{huber1964robust}.
In practice we use $\delta = 0.1 \cdot \max|\tensor{R}_0|$
(Section~\ref{subsec:gradient}) as a scale-adaptive initialisation.

% ------------------------------------------------------------
\subsection{Geometric Interpretation}
\label{subsec:norm_geometry}
% ------------------------------------------------------------

Each norm induces a different geometry on the residual space
$\R^{I_1 \times \cdots \times I_N}$.

\begin{itemize}
  \item \textbf{Frobenius ($\ell_2$)}: unit ball is a hypersphere.
        The projection onto the low-rank manifold $\mathcal{M}_r$ is
        the Euclidean nearest point, computed exactly by HOSVD/HOOI.

  \item \textbf{$\ell_1$}: unit ball is a cross-polytope.
        The solution places large residuals at the vertices of the
        polytope, concentrating error on a small number of entries
        (sparse residual structure).

  \item \textbf{Huber}: unit ball interpolates between the two.
        The transition radius $\delta$ determines how quickly the ball
        transitions from a hypersphere centre to a polytope surface.
\end{itemize}

\noindent
This geometric picture predicts a key experimental observation: under
sparse corruption, $\mathcal{L}_1$-Tucker concentrates its residual
on the corrupted entries, leaving the clean entries nearly exactly
reconstructed, whereas $\mathcal{L}_F$-Tucker spreads the residual
evenly across all entries.
Section~\ref{sec:experiments} verifies this prediction via error maps.

% ------------------------------------------------------------
\subsection{Norm Comparison Under Structured Noise}
\label{subsec:norm_comparison}
% ------------------------------------------------------------

Let the observed tensor be
$\tensor{X} = \tensor{X}^* + \tensor{S} + \tensor{G}$,
where $\tensor{X}^*$ is the clean low-rank signal,
$\tensor{S}$ is a sparse corruption tensor with
$\|\tensor{S}\|_0 = s$ nonzeros,
and $\tensor{G}$ is dense Gaussian noise with variance $\sigma^2$.

Table~\ref{tab:norm_noise} summarises the expected behaviour of each
estimator under this model.

\begin{table}[h]
  \centering
  \caption{Expected reconstruction behaviour under structured noise.}
  \label{tab:norm_noise}
  \begin{tabular}{lccc}
    \toprule
    & \multicolumn{3}{c}{Noise type} \\
    \cmidrule(lr){2-4}
    Loss & Dense Gaussian & Sparse outliers & Both \\
    \midrule
    $\mathcal{L}_F$ & Optimal (MLE) & Poor & Poor \\
    $\mathcal{L}_1$ & Suboptimal    & Robust ($\varepsilon^*{=}0.5$) & Moderate \\
    $\mathcal{L}_H$ & Near-optimal  & Robust ($\delta$-dependent) & Good \\
    \bottomrule
  \end{tabular}
\end{table}

\noindent
The key prediction is that the relative advantage of
$\mathcal{L}_1$ and $\mathcal{L}_H$ over $\mathcal{L}_F$ should grow
monotonically with the fraction of corrupted entries $s / \prod_n I_n$.
This is tested directly in the noise robustness experiment
(Section~\ref{subsec:exp_noise}).

% % ============================================================
% %  sections/chapter_04_ko.tex
% %  4장: 노름 분석
% % ============================================================

% \section{노름 분석}
% \label{sec:norms}

% 재구성 손실 $\mathcal{L}$의 선택은 단순한 계산상의 편의가 아니라,
% 잡음 분포에 대한 명시적 가정을 담고 있으며
% 근사가 어떤 오차를 우선적으로 최소화할지를 결정한다.
% 본 장에서는 \ref{subsec:losses}절에서 소개한 세 가지 손실 함수를
% 통계적·기하학적·수치적 관점에서 분석하고,
% 각 노름 하에서 Tucker 재구성에 미치는 영향을 유도한다.

% % ------------------------------------------------------------
% \subsection{프로베니우스 노름: 통계적·기하학적 특성}
% \label{subsec:frob_analysis}
% % ------------------------------------------------------------

% \paragraph{최대 우도 해석.}
% 관측 텐서가
% $\tensor{X} = \hat{\tensor{X}} + \tensor{E}$이고
% $e_{i_1\cdots i_N} \overset{\text{i.i.d.}}{\sim} \mathcal{N}(0, \sigma^2)$라
% 가정하면, 로그 우도는
% \begin{equation}
%   \log p(\tensor{X} \mid \hat{\tensor{X}})
%   \;=\; -\frac{1}{2\sigma^2} \frobn{\tensor{X} - \hat{\tensor{X}}}^2
%         + \text{상수}
%   \label{eq:gauss_ll}
% \end{equation}
% 이므로, $\hat{\tensor{X}}$의 최대 우도 추정량은 프로베니우스 최소화와
% 정확히 일치한다.
% 이 잡음 모델 하에서 HOOI는 최적이다.

% \paragraph{유니터리 불변성과 ALS 단순화.}
% 유니터리 불변성~\eqref{eq:unitary_inv}에 의해
% 인수 행렬이 정규직교일 때 목적 함수는 다음과 같이 분해된다:
% \begin{equation}
%   \frobn{\tensor{X} - \hat{\tensor{X}}}^2
%   \;=\; \frobn{\tensor{X}}^2 - \frobn{\tensor{G}}^2.
%   \label{eq:frob_decomp}
% \end{equation}
% 즉, \emph{재구성 오차를 최소화하는 것은 코어 텐서가 포착하는
% 에너지를 최대화하는 것과 동치}이다.
% 이것이 HOSVD와 HOOI가 활용하는 변분적 특성화이다.

% \paragraph{이상치 민감도.}
% 제곱 손실은 각 잔차 $r = x - \hat{x}$에 $|r|$에 비례하는 가중치를
% 부여하므로, 단 하나의 큰 이상치가 목적 함수를 지배하여
% 모든 인수 행렬을 왜곡할 수 있다.
% 형식적으로, 프로베니우스 추정량의 \emph{영향 함수(influence function)}는
% 비유계(unbounded)이다~\cite{hampel1986robust}:
% \begin{equation}
%   \mathrm{IF}(r;\, \mathcal{L}_F) \;=\; r.
%   \label{eq:if_frob}
% \end{equation}
% 이는 총오차(gross error)에 대한 강인성이 전혀 없음을 확인시켜 준다.

% % ------------------------------------------------------------
% \subsection{\texorpdfstring{$\ell_1$}{L1} 노름: 강인성과 비매끄러움}
% \label{subsec:l1_analysis}
% % ------------------------------------------------------------

% \paragraph{최대 우도 해석.}
% $e_{i_1\cdots i_N} \overset{\text{i.i.d.}}{\sim} \mathrm{Laplace}(0, b)$이면
% 로그 우도는
% \begin{equation}
%   \log p(\tensor{X} \mid \hat{\tensor{X}})
%   \;=\; -\frac{1}{b}\onenorm{\tensor{X} - \hat{\tensor{X}}}
%         + \text{상수}
%   \label{eq:laplace_ll}
% \end{equation}
% 이므로, $\ell_1$ 최소화는 라플라스 잡음 하의 최대 우도 추정과 동치이다.
% 라플라스 분포의 가우시안 대비 두꺼운 꼬리(heavy tail)는
% $\mathcal{L}_1$에 내재된 이상치 허용 능력을 반영한다.

% \paragraph{분해점 (Breakdown point).}
% 추정량의 \emph{분해점} $\varepsilon^*$은 추정값이 임의로 편향되기 전까지
% 허용할 수 있는 임의 손상 관측값의 비율이다.
% 프로베니우스 추정량은 $\varepsilon^* = 0$ (단 하나의 이상치로
% 추정값이 무한히 멀어질 수 있음)인 반면,
% $\ell_1$ 추정량은 위치 문제에서 $\varepsilon^* = 0.5$를
% 달성한다~\cite{maronna2019robust}.
% Tucker 설정에서 이는 텐서 원소의 최대 절반까지의 희소 손상에 대한
% 강인성으로 해석된다.

% \paragraph{부분미분과 최적성 조건.}
% $\mathcal{L}_1$은 비매끄러우므로 고전적 경사 조건이 적용되지 않는다.
% $|r|$의 $r$에서의 부분미분(subdifferential)은
% \begin{equation}
%   \partial|r|
%   \;=\; \begin{cases}
%           \{\sign(r)\} & r \neq 0, \\
%           [-1,\, 1]    & r = 0
%         \end{cases}
%   \label{eq:subgrad_l1}
% \end{equation}
% 이며, $\mathcal{L} = \mathcal{L}_1$인~\eqref{eq:general_opt}의 최적성 조건은
% 각 인수에 대한 목적 함수의 부분미분에 0이 포함되어야 함을 요구한다.
% 이 조건이 수렴 시 IRLS (\ref{subsec:irls}절)가 만족하고자 하는 조건이다.

% \paragraph{영향 함수.}
% $\ell_1$ 추정량의 영향 함수는 유계(bounded)이다:
% \begin{equation}
%   \mathrm{IF}(r;\, \mathcal{L}_1) \;=\; \sign(r).
%   \label{eq:if_l1}
% \end{equation}
% 이는 단일 이상치가 추정값을 진실의 유한 이웃 밖으로 밀어낼 수 없음을 확인한다.

% % ------------------------------------------------------------
% \subsection{Huber 손실: 매끄러운 보간}
% \label{subsec:huber_analysis}
% % ------------------------------------------------------------

% \paragraph{정의와 성질.}
% Huber 손실 $\rho_\delta$~\eqref{eq:huber}는 다음 성질을 만족한다:
% \begin{itemize}
%   \item $\rho_\delta \in C^1(\R)$: 모든 점에서 연속 미분 가능하므로
%         부분기울기 없이 경사 하강법을 적용할 수 있다.
%   \item $\rho_\delta(r) = \tfrac{1}{2}r^2$ ($|r| \leq \delta$):
%         원점 근방에서 이차 함수로 가우시안 소잡음과 일치한다.
%   \item $\rho_\delta(r) \sim \delta|r|$ ($|r| \gg \delta$):
%         꼬리에서 선형으로 $\ell_1$의 강인성을 계승한다.
%   \item $\delta \to \infty$이면 $\rho_\delta \to \tfrac{1}{2}|\cdot|^2$,
%         $\delta \to 0$이면 $\rho_\delta \to \delta|\cdot|$:
%         Huber는 순수 프로베니우스와 순수 $\ell_1$ 사이를 보간한다.
% \end{itemize}

% \paragraph{영향 함수.}
% \begin{equation}
%   \mathrm{IF}(r;\, \mathcal{L}_H)
%   \;=\; \psi_\delta(r)
%   \;=\; \begin{cases}
%           r                    & |r| \leq \delta\text{인 경우,} \\
%           \delta\,\sign(r) & |r| > \delta\text{인 경우}
%         \end{cases}
%   \label{eq:if_huber}
% \end{equation}
% 로 유계(프로베니우스와 달리)이면서 매끄럽다($\ell_1$과 달리).

% \paragraph{임계값 선택.}
% $\delta$는 이차 구간과 선형 구간의 경계를 제어한다.
% 데이터 적응적 선택은
% \begin{equation}
%   \delta^*
%   \;=\; c \cdot \widehat{\mathrm{MAD}}(\tensor{R}),
%   \qquad
%   \widehat{\mathrm{MAD}} = \mathrm{median}\!\left(|r_{i_1\cdots i_N}|
%   \right)
%   \label{eq:delta_mad}
% \end{equation}
% 이며, 가우시안 잡음에서 $c \approx 1.345$가 표준 효율 상수이다~\cite{huber1964robust}.
% 본 실험에서는 스케일 적응적 초기화로
% $\delta = 0.1 \cdot \max|\tensor{R}_0|$를 사용한다
% (\ref{subsec:gradient}절).

% % ------------------------------------------------------------
% \subsection{기하학적 해석}
% \label{subsec:norm_geometry}
% % ------------------------------------------------------------

% 각 노름은 잔차 공간
% $\R^{I_1 \times \cdots \times I_N}$에 서로 다른 기하 구조를 부여한다.

% \begin{itemize}
%   \item \textbf{프로베니우스 ($\ell_2$)}: 단위 볼이 초구(hypersphere)이다.
%         저랭크 다양체 $\mathcal{M}_r$로의 투영이 유클리드 최근접점이 되어
%         HOSVD/HOOI로 정확하게 계산된다.

%   \item \textbf{$\ell_1$}: 단위 볼이 교차 다면체(cross-polytope)이다.
%         해는 큰 잔차를 다면체의 꼭짓점에 배치하여
%         소수의 원소에 오차를 집중시킨다 (희소 잔차 구조).

%   \item \textbf{Huber}: 단위 볼이 두 형태의 중간을 보간한다.
%         전환 반경 $\delta$가 볼의 중심부 초구에서 표면 다면체로의
%         전환 속도를 결정한다.
% \end{itemize}

% 이 기하학적 관점은 핵심 실험 관찰을 예측한다:
% 희소 손상 하에서 $\mathcal{L}_1$-Tucker는 잔차를 손상된 원소에
% 집중시켜 깨끗한 원소를 거의 정확하게 재구성하는 반면,
% $\mathcal{L}_F$-Tucker는 잔차를 모든 원소에 균등하게 분산시킨다.
% \ref{sec:experiments}장에서 오차 지도(error map)를 통해 이 예측을 검증한다.

% % ------------------------------------------------------------
% \subsection{구조적 잡음 하에서의 노름 비교}
% \label{subsec:norm_comparison}
% % ------------------------------------------------------------

% 관측 텐서를
% $\tensor{X} = \tensor{X}^* + \tensor{S} + \tensor{G}$로 모델링한다.
% 여기서 $\tensor{X}^*$는 깨끗한 저랭크 신호,
% $\tensor{S}$는 비영 원소가 $\|\tensor{S}\|_0 = s$개인 희소 손상 텐서,
% $\tensor{G}$는 분산 $\sigma^2$의 밀집 가우시안 잡음이다.

% 표~\ref{tab:norm_noise}는 이 모델 하에서 각 추정량의 예상 거동을 요약한다.

% \begin{table}[h]
%   \centering
%   \caption{구조적 잡음 하에서의 예상 재구성 거동.}
%   \label{tab:norm_noise}
%   \begin{tabular}{lccc}
%     \toprule
%     & \multicolumn{3}{c}{잡음 유형} \\
%     \cmidrule(lr){2-4}
%     손실 함수 & 밀집 가우시안 & 희소 이상치 & 혼합 \\
%     \midrule
%     $\mathcal{L}_F$ & 최적 (MLE) & 불량 & 불량 \\
%     $\mathcal{L}_1$ & 준최적      & 강인 ($\varepsilon^*{=}0.5$) & 보통 \\
%     $\mathcal{L}_H$ & 준최적      & 강인 ($\delta$ 의존) & 양호 \\
%     \bottomrule
%   \end{tabular}
% \end{table}

% 핵심 예측은 $\mathcal{L}_1$과 $\mathcal{L}_H$의 $\mathcal{L}_F$ 대비
% 상대적 이점이 손상 원소 비율 $s / \prod_n I_n$에 따라
% 단조 증가해야 한다는 것이다.
% 이는 \ref{subsec:exp_noise}절의 잡음 강인성 실험에서 직접 검증한다.